\chapter{The viewers in Medorama}
\label{chap:viewers-in-medorama}

This chapter describes how Medorama opens its viewers and what it does around them: the
patient's context, the gateway that stands between users and VitalStream, and Privacy Mode. It is
written for users and administrators. Each viewer's own features are described in
\cref{chap:vitalstream}, \cref{chap:ecg-viewer} and \cref{chap:pacs-viewer}.

Medorama shows recordings, ECGs and images in its own viewers, as tabs of its page:
\ui{VitalStream} for waveform recordings, \ui{ECG} for 12-lead ECGs, and \ui{PACS} for DICOM
studies, pictures, videos and PDFs. They are the same programs that run alone
(\cref{part:viewers}), unchanged. Medorama opens them at the file to show, behind its own sign-in
and its privacy gate, and tells them whose file it is.

\begin{table}[htbp]
\caption{The viewers, inside Medorama.}
\label{tab:viewers-in-medorama-viewers}
\begin{tabularx}{\linewidth}{@{}P{2.0cm}LP{4.6cm}@{}}
\toprule
{\tabhead Viewer} & {\tabhead Shows} & {\tabhead Opened from} \\
\midrule
\ui{VitalStream} & Waveform recordings (ICU monitoring, Holter, EEG, sleep studies): every
  channel, vital signs, alarms or annotations and breaks on one time axis, from days down to
  single samples & \ui{View} by a recording; a recording's entry on the patient page, at its
  moment; a database's row; the tab itself \\
\ui{ECG} & 12-lead ECGs on standard paper (3 $\times$ 4 with a rhythm strip, 6 $\times$ 2,
  12 $\times$ 1), with calipers, and their report & \ui{Open in ECG} on the patient page; a
  database's row \\
\ui{PACS} & DICOM studies (echo and angiography loops at their frame rate, CT, MR, X-ray,
  reports), videos, pictures and PDFs, with calipers in millimeters & \ui{Open in PACS} on the
  patient page; a database's row \\
\bottomrule
\end{tabularx}
\end{table}

\section{Opening a file in a viewer}
\label{sec:viewers-in-medorama-opening}

A file opens in its viewer from any of the places where Medorama shows it:
\begin{itemize}
\item \textbf{the patient's page} (\cref{chap:patients}): \ui{Open in ECG} and \ui{Open in PACS}
  in the detail of an ECG, an echo loop, an image or a series; for a recording, \ui{Open in
  VitalStream} (at its start) or a click on its overview (at that moment);
\item \textbf{the Databases tab} (\cref{chap:databases-tab}): the link by a row's file
  (\ui{VitalStream}, \ui{ECG}, \ui{PACS});
\item \textbf{an encounter's panel} and the recordings of \ui{Import}: \ui{View} by a recording,
  which opens it whole.
\end{itemize}

The viewer's tab then shows the file. Each viewer keeps its state while the user goes to other
screens, and shows another file without reloading when it is sent one. The \ui{ECG}
and \ui{PACS} tabs are dimmed until something is opened in them (``Opens what a patient's page
sends here''); until then, they say ``Open an ECG, an echo or an image from a patient's page.''
The \ui{VitalStream} tab can be opened directly: it shows the viewer's own page, with its
\ui{Open record\dots} and \ui{Open multiple records\dots} dialogs, limited by the gateway's
rules (\cref{sec:viewers-in-medorama-gateway}).

\paragraph{Which viewer shows which file.} The server decides, from the item's category
(\cref{chap:patients}) and the file's kind:
\begin{itemize}
\item the recordings of the ICU, Holter, EEG and Sleep categories open in VitalStream;
\item the files of the ECG category (WFDB records, GE MUSE XML, DICOM ECGs, EDF and BDF files,
  JSON, pictures and PDFs) open in the ECG viewer;
\item the files of the Echo, Cath lab and Imaging categories (DICOM files and folders, videos,
  pictures and PDFs) open in the PACS viewer;
\item any other file is opened by the browser, as it is (\ui{Open}).
\end{itemize}
The detail of the patient's page also shows most files itself: a 12-lead ECG drawn on paper by
the ECG viewer and framed in the detail, or a picture, a PDF or a video.

\paragraph{Addresses.} Each viewer tab has its own address (\cref{chap:page}):
\begin{itemize}
\item \code{\#/viewer?path=\dots\&t0=\dots\&t1=\dots}: a recording at a time span (Unix
  seconds);
\item \code{\#/ecg?src=\dots}: an ECG file; \code{\#/pacs?src=\dots}: one file;
  \code{\#/pacs?list=\dots}: a folder, its files listed by the API;
\item and, when the file was opened from a patient's page or a database's row,
  \code{patient=\dots\&entry=\dots}: whose file it is (\cref{sec:viewers-in-medorama-context}).
\end{itemize}
Back and forward move between the patient's page and the viewers.

\section{Whose file: the host's context}
\label{sec:viewers-in-medorama-context}

Inside Medorama, a file's name is of little use. Instead, the viewers' title says whose file it
is, as clinical systems show a banner to prevent wrong-patient errors. The title has three
parts:
\begin{itemize}
\item \textbf{who}: the patient as the user may see them (the name, the patient ID, the MRN,
  sex and age for users with the PHI permission out of Privacy Mode; in Privacy Mode, the
  placeholders and the pseudonym), never cut short;
\item \textbf{what}: the entry's title, its date and time, and its encounter;
\item \textbf{where among the patient's files}: ``3 of 4'', the place of this file among the
  patient's files of that kind, in time order.
\end{itemize}

\begin{description}
\item[\ui{$\langle$ Patient}] returns to the entry on the patient's page, as it was.
\item[\ui{$\langle$} \ui{$\rangle$}] step to the patient's previous and next file of that kind
  (their recordings in VitalStream; their ECGs, or their images, in the other viewers).
\end{description}

Medorama passes this context in the viewer's address (\code{who}, \code{what}, \code{pos},
\code{nav}; \cref{chap:viewer-api}). When one of these buttons is used, the viewer only reports
it to the page, by a message to its host of the same origin, and the page navigates. The viewer
never sees patient IDs or addresses. A viewer opened without a patient (the VitalStream tab
opened directly) shows the file's name.

\paragraph{Themes.} Medorama's page and every viewer have a light and a dark theme. The three
viewers take the page's theme from their address (\code{theme=light} or \code{dark},
\cref{sec:viewer-api-addresses}). Medorama passes its page's theme (\cref{chap:page}) to every
viewer tab, VitalStream's included, when the tab opens and whenever the theme changes. A viewer
applies it without remembering it as its own choice, and VitalStream's own theme button still
switches it.

\section{VitalStream in Medorama}
\label{sec:viewers-in-medorama-vitalstream}

Medorama runs one VitalStream viewer for its users and shows it in its \ui{VitalStream} tab.
The setting \code{viewer} chooses which one: the C++, Rust or Python implementation, a program's
path, none, or a viewer running elsewhere. By default it is the one the installer built
(\cref{chap:administration}).

\begin{itemize}
\item \textbf{It starts at first use}, at the first \ui{View} or the first visit of the tab
  (``Starting the viewer\dots''), as its own process on 127.0.0.1 and a free port. It is
  restarted if it stops, and stopped with the catalog. If it cannot start, the tab says why
  (``The viewer cannot start: \dots'').
\item \textbf{Its settings} are a file that Medorama writes at each start
  (\file{<home>/viewer/vitalstream.toml}). It names the catalog's index folder (so that every
  viewer writes and reads the same index), its adapter folder, its time zone and
  \code{data\_roots}, and requires a password, so that a viewer built without password support
  refuses to run under the catalog. A user's own \file{vitalstream.toml} cannot conflict with
  these settings. The limits on a folder of records (\cref{sec:vitalstream-limits}) are
  VitalStream settings; Medorama has no such settings, and indexes a recording's folder within
  the defaults.
\item \textbf{Its password} has a signing key that only the catalog knows. The catalog signs
  every request it passes on, users never see the viewer's sign-in, and the viewer's port is of
  no use to anyone else on the machine.
\item \textbf{No viewer.} When none is installed beside the catalog, or the setting is
  \code{none}, the \ui{VitalStream} tab and the \ui{View} buttons are grayed out, with the
  reason in their tooltip (``No viewer is linked to the catalog'').
\item \textbf{A viewer elsewhere} (the setting \code{viewer\_url}). The catalog runs no viewer,
  the tab is absent, and \ui{View} opens that viewer in a browser tab of its own.
\end{itemize}

\section{The gateway}
\label{sec:viewers-in-medorama-gateway}

The VitalStream viewer is reached only through Medorama, at \file{/viewer/}, so that users have
one address and one sign-in, and browsing is controlled. After the catalog's sign-in, Medorama's
\emph{gateway} decides each request by a rule (\cref{tab:viewers-in-medorama-gateway}). Anything
else, the viewer's own \code{login} and \code{logout} included, is refused (default deny); so is
any endpoint that a newer viewer may add, until it is given a rule. The gateway's place in the
architecture is described in \cref{sec:architecture-gateway}.

\begin{xltabular}{\linewidth}{@{}P{4.6cm}P{3.6cm}L@{}}
\caption{The requests the gateway passes to VitalStream.}
\label{tab:viewers-in-medorama-gateway}\\
\toprule
{\tabhead Viewer request} & {\tabhead Admins} & {\tabhead Analysts, viewers} \\
\midrule
\endfirsthead
\toprule
{\tabhead Viewer request} & {\tabhead Admins} & {\tabhead Analysts, viewers} \\
\midrule
\endhead
\bottomrule
\endfoot
The page and its files & Passed & Passed \\
\code{GET api/v1/config} & As the viewer answers & Its places, start and home folders point
  into the folders of the catalog's recordings \\
\code{GET api/v1/browse}, \code{complete} (the Open dialogs) & Inside \code{data\_roots}
  (anywhere when it is empty) & Only folders that lead to recordings in the catalog, listing
  only those recordings \\
\code{GET api/v1/scan} (what a folder of records holds, before it opens) & Inside
  \code{data\_roots} & A recording in the catalog; refused in Privacy Mode \\
\code{POST api/v1/open} & A recording inside \code{data\_roots} & A recording in the
  catalog \\
\code{POST api/v1/cancel} (stop an index being built) & Datasets the user opened through the
  gateway & The same \\
\code{GET api/v1/datasets/\{key\}\dots} (windows, vital signs, markers) & Datasets the user
  opened through the gateway & The same \\
Anything else & Refused & Refused \\
\end{xltabular}

\begin{itemize}
\item Every recording opened is written to the audit log (\code{view}).
\item Paths are compared after resolving symbolic links, as the viewers do.
\item The viewer's page may be framed only by the catalog's own page.
\item The viewer receives paths and recordings, never identifiers.
\item A request body over 1\,MiB is answered with \code{413}. A \code{POST api/v1/open} whose
  body is not sent as \code{application/json} is answered with \code{415}, and one whose body is
  not valid JSON with \code{400}.
\end{itemize}

The ECG and PACS viewers need no gateway. They are static pages that Medorama serves (at
\file{/viewers/ecg/} and \file{/viewers/pacs/}), and they read only the files of an entry,
through \code{api/v1/patients/\{id\}/file/\{item\}}. Medorama checks these requests like every
other request: only the patient's file, inside the data roots, and, for a WFDB header, only the
signal files it names, are served (\cref{chap:http-api}).

\section{Privacy Mode in the viewers}
\label{sec:viewers-in-medorama-privacy}

\begin{phibox}
The viewers show no patient identity of their own: whatever a file holds, they never display a
name or an identifier from it. Inside Medorama, the identity they show is the host's context
(\cref{sec:viewers-in-medorama-context}), as the server sent it for the user. In Privacy Mode,
the server also de-identifies what they receive:
\begin{itemize}
\item \textbf{Files.} WFDB headers, GE MUSE XML and EDF headers are de-identified, and the texts
  of EDF+ annotations removed. DICOM headers are rewritten following the PS3.15 Basic Profile,
  with the pixels untouched. File names are replaced by tokens (\code{mo-f0003.dcm}). Pictures,
  videos and PDFs are sent as they are, and anything else is refused.
\item \textbf{An amber warning triangle} by the \ui{ECG} or \ui{PACS} tab marks a file whose
  pixels or pages are sent as they are (a picture, a PDF, a video or a DICOM image). Such a file
  may show burned-in text, which Privacy Mode cannot hide.
\item \textbf{VitalStream.} Recordings are opened by token (\code{mo-rec:12}), never by path.
  Browsing, completing and scanning folders are refused, because folder and file names may
  identify. The times the page sends, which it knows moved back, are moved forward again before
  they reach the viewer, and the times in the viewer's answers are moved back. A recording's path
  and name are replaced. The texts of EDF+ annotations are treated like the text of notes.
\end{itemize}
Every answer then passes the gate's last layer like any other. The rules are given in full in
\cref{sec:protecting-files} and \cref{sec:protecting-vitalstream}, and the design and its limits
in \cref{chap:protecting} and \cref{chap:security}.
\end{phibox}

\section{When the data change}
\label{sec:viewers-in-medorama-changes}

When a database is deleted, in this browser or in another, every open page learns of it within
seconds (\cref{chap:page}). An ECG or PACS tab whose file is gone closes and says so (``What
this ECG tab showed was removed from the catalog (its database was deleted).''), and the tab is
dimmed again.
